\section{Lessons: Wiring Is the Trusted Base}\label{sec:lessons}
The theorems of Section~\ref{sec:mech} speak about what the certified
verifier checks; what reaches it, and about which values, is decided by
unverified code, and every discrepancy we have found between a certified
pipeline and a deployed one, in Belenios (Section~\ref{sec:belenios}) and
in our own earlier verifiers, lay in that code rather than in a
verification equation. Three parts of the boundary carry the security of
the Schulze scheme.

\paragraph{The decoder}
Ballot soundness is a theorem about elements of the group type. A decoder
that admits a ciphertext component outside the order-$q$ subgroup, as
Belenios's shuffle verifier does, hands the certified code an input about
which no theorem speaks. The Schulze verifier's decoder therefore checks
the order of every group element and the range of every field element
before anything reaches the certified functions; this check is the first
candidate for verification, since a decoder is a parser with a simple
specification.

\paragraph{The wiring}
Two decisions outside the theorems carry the security of the scheme. The
verifier must \emph{derive} the $\delta$ ciphertexts from the position
ciphertexts and never accept them as inputs (Section~\ref{sec:ballot}); a
driver that read $\mathrm{Enc}(\delta_{ij})$ from the ballot would verify
every proof and count a ballot twice. And the Fiat--Shamir challenge of a
ballot must bind the voter's identity, or a ballot can be replayed under
another credential; the mechanised challenge covers the ciphertexts and
announcements, and the identity is the driver's responsibility
\TODO{implement and test}. Both are tested the way we tested Belenios:
the certified prover, which takes its witness and randomness as inputs,
produces the double-counting ballot and the replayed ballot, and the
extracted verifier must reject them. Wiring is protocol logic that no
current formalisation covers end to end; the \texttt{count} judgement of
Section~\ref{sec:mech} covers the aggregation and the final check, and
extending it to the derivation of the statements each proof is about is
where we would put the next proof.

\paragraph{The challenge closures}
The non-interactive theorems take the hash as a parameter. A closure that
hashes the wrong announcement is invisible to a verifier that trusts the
stored challenges and visible only to one that recomputes them, which the
Schulze verifier does; proving that the client hashes exactly the
announcements it sends (\texttt{nizk\_ballot\_valid} is stated with the
challenge equal to the hash of the transcript) is what makes the
recomputation sound.
